% Pursuing Stacks (cotes 134-2 to 134-8): archive cote 134-7, site page 93 (PDF page 94), typescript page 535.
% typescript. Draft compared: ps6.tex:4674-4736.
% AI-transcribed by Claude Opus 5.5 (claude-opus-5-5, adaptive thinking, effort high), Message Batches API, batch ps4 (msgbatch_012Aoc1NuskATPsyBeKSYK7J), 2026-11-12; not yet checked by eye.
\scanpage{134-7/93}
\mspage{535}

latter by either L\hand{$\pi$}. , or LH. . Now the point is that the properties
of this subcategory, with respect to the inner structure of Hotab$_o$(k)
(involving the``left'' and the ``right'' homotopy operations, notably the
suspension and the loop functors\over{x}{\unsure{,}} respectively) are by no means autosym-%
metric. The main dissymetry, it would seem, turns out in this, that
(\over{5}{4}) is stable under the loop functor, and by no means under the suspen-%
sion functor. This \xout{implies that}\inserted[typed]{is the reason why,} even if we assume that (in analogy
to what happens for the categories (3)) the inclusion functors from the
categories (5) into Hot$_o$(k) do admit the relevant (namely right) adjoint
(which I greatly doubt anyhow ...), and thus give rise for any object X
to an increasing filtration (*) , the Postnikov-type relationships\xout{\unsure{betw}}
between X$_n$ and X$_{n+1}$ cannot be quite dualized to a similar relationship
between $_n$X and $_{(n+1)}$X . We may think \inserted{next,} of course\hand{,} of defining an increa-%
sing sequence of subcategories \xout{of} (5) of H\over{o}{\unsure{o}}t$_o$(k) in terms of LH. and a
correspondance sequence of subcategories of Hotab$_o$(k), different from
the categories (\over{5}{4}), and stable under suspensions. \xout{Here the natural}
\xout{choice would be to} But there \unsure{doesn't|does'nt} seem to be anything reasonable \xout{aro}
around along these lines.

To sum up, it does'nt seem one should overemphasize the somewhat
startling symmetry which appeared in section 124 between ``homotopy''
\xout{invariants} embodied in L\hand{$\pi$}. , and``homology''\xout{invariants} embodied in LH. -
in some essential respects, it would seem that the corresponding two
functors do have non mutually symmetric properties. I guess I have to
apologize for having taken that long for coming to a conclusion which,
presumably, must be felt as a kind of self-evidence \over{\illegible{1}}{by} \inserted{all} homotopy
people ! \margin{(line drawn under this line and bracket at the left of the next line, separating the sections)}

127. Before leaving (for the time being) the topic of schematic
homotopy types and schematization, I would like still to add a few com-%
\xout{ments} ments, about various possibilities for working with different kind
of models for defining schematic homotopy types. My point here is not
in \xout{changing}\inserted[typed]{replacing} the \inserted[typed]{basic} test-category we are working with, here \hand{$\Delta$} , by some
other (say the category \hand{\unsure{$\bigcirc$}} of \inserted[typed]{standard} hemispheres) - \xout{\unsure{we'll keep}} this choice I
feel should be more or less irrelevant, so we may as well keep \hand{$\Delta$} .
Thus we are going to work with semisimplicial objects, and the \deleted{whole} \inserted{main}
question then is (for a given ground ring k) to say \xout{wh} precisely what
kind of objects we are allowing (or imposing !) as components for our
complexes. At any rate they should be ``objects over k'', and the most

\margin{93}
